\documentclass{article}
\usepackage{amsmath, amssymb, bm}

\begin{document}

\title{Extracted Equations from PDF}
\author{}
\date{}
\maketitle

\section{Trigonometry}
\begin{align}
    c^2 &= a^2 + b^2 - 2ab \cos \gamma \\
    \cos(\alpha + \beta) &= \cos\alpha \cos\beta - \sin\alpha \sin\beta \\
    \sin(\alpha + \beta) &= \cos\alpha \sin\beta + \sin\alpha \cos\beta
\end{align}

\section{Damped Harmonic Oscillator}
\begin{equation}
    \ddot{x} + 2\gamma\dot{x} + \omega_0^2 x = f(x)
\end{equation}

\section{Fourier Transform Solution}
\begin{equation}
    x_p = \frac{1}{\sqrt{2\pi}} \int_{-\infty}^{\infty} \frac{\hat{f}(\omega) e^{i\omega t} d\omega}{\omega_0^2 - \omega^2 + 2i\gamma\omega}
\end{equation}

\section{Forced Oscillation}
\begin{align}
    x_p &= \frac{f_0 \cos(\omega_d t + \varphi)}{\sqrt{(\omega_0^2 - \omega_d^2)^2 + (2\gamma\omega_d)^2}} \\
    \varphi &= \tan^{-1} \left( \frac{2\gamma\omega_d}{\omega_0^2 - \omega_d^2} \right)
\end{align}

\section{Matrix Inverse}
\begin{equation}
    A = \begin{pmatrix} a & b \\ c & d \end{pmatrix} \Rightarrow A^{-1} = \frac{1}{\det A} \begin{pmatrix} d & -b \\ -c & a \end{pmatrix}
\end{equation}

\section{Levi-Civita Symbol}
\begin{equation}
    (v \times u)_k = \epsilon_{ijk} v_i u_j
\end{equation}

\section{Legendre Transform}
\begin{align}
    \frac{df}{dx} &= f'(x) \equiv s(x) \\
    g(s) &= -x(s) \cdot s + f(x(s)) \\
    \frac{\partial g}{\partial s} &= -x(s) \\
    V(x) &= g(s(x)) + xu(x)
\end{align}

\section{Euler's Theorem}
\begin{align}
    x \cdot \nabla f(x) &= k f(x) \\
    x f(x) &= k \frac{\partial f}{\partial x}
\end{align}

\section{Elliptic Area Integral}
\begin{equation}
    \int_0^{2\pi} \int_0^1 dA = \int_0^{2\pi} \int_0^1 abr \, dr \, d\theta
\end{equation}

\section{Rotation Matrices}
\begin{equation}
    R_z(\phi) =
    \begin{pmatrix}
        \cos\phi & \sin\phi & 0 \\
        -\sin\phi & \cos\phi & 0 \\
        0 & 0 & 1
    \end{pmatrix}
\end{equation}
\mathbf{\Omega} =
\begin{pmatrix}
\omega_1 \\
\omega_2 \\
\omega_3
\end{pmatrix}
=
\begin{pmatrix}
\dot{\phi} \sin\theta \sin\psi + \dot{\theta} \cos\psi \\
\dot{\phi} \sin\theta \cos\psi + \dot{\theta} \sin\psi \\
\dot{\psi} + \dot{\phi} \cos\theta
\end{pmatrix}

\section{Poisson Brackets}
\begin{equation}
    \{f, g\}_{pq} := \sum \left( \frac{\partial f}{\partial p_i} \frac{\partial g}{\partial q_i} - \frac{\partial f}{\partial q_i} \frac{\partial g}{\partial p_i} \right)
\end{equation}

\section{Angular Momentum}
\begin{align}
    L &= r \times p \\
    \tau &= r \times F
\end{align}

\section{Lagrange’s Equations}
\begin{align}
    L &= K - V \\
    \frac{d}{dt} \left( \frac{\partial L}{\partial \dot{q}_i} \right) - \frac{\partial L}{\partial q_i} &= 0
\end{align}

\section{Energy Conservation}
\begin{align}
    E &= \sum q_i \frac{\partial L}{\partial \dot{q}_i} - L \\
    \frac{\partial L}{\partial t} = 0 &\Rightarrow \frac{dE}{dt} = 0
\end{align}

\section{Time Integral for Motion}
\begin{align}
    t &= \int_{x_0}^{x} \frac{dx}{\sqrt{\frac{2}{m} (E - V(x))}} \\
    \tau &= 2 \int_{x_1(E)}^{x_2(E)} \frac{dx}{\sqrt{\frac{2}{m} (E - V(x))}}
\end{align}

\section{Hamilton’s Equations}
\begin{align}
    H &= \sum p_j \dot{q}_j - L \\
    \dot{q}_i &= \frac{\partial H}{\partial p_i} \\
    \dot{p}_i &= -\frac{\partial H}{\partial q_i}
\end{align}

\section{Hamilton-Jacobi Equation}
\begin{equation}
    H(q_j, \frac{\partial S}{\partial q_j}, t) + \frac{\partial S}{\partial t} = 0
\end{equation}

\section{Euler’s Equations for Rotation}
\begin{align}
    I_1 \dot{\omega}_1 + (I_3 - I_2) \omega_2 \omega_3 &= \tau_1 \\
    I_2 \dot{\omega}_2 + (I_1 - I_3) \omega_3 \omega_1 &= \tau_2 \\
    I_3 \dot{\omega}_3 + (I_2 - I_1) \omega_1 \omega_2 &= \tau_3
\end{align}

\section{Virial Theorem}
\begin{equation}
    2 \langle K \rangle = k \langle V \rangle
\end{equation}

\section{Rigid Body Inertia Moments}

\textbf{Mathieu (Parametric) Oscillator}
\[
\ddot{x} \;+\;\bigl[\omega_0^2 + \varepsilon \cos(\gamma t)\bigr]\,x \;=\; 0.
\]
For small \(\varepsilon\), one can analyze solutions via expansions in \(\sin(\omega_0 t+\delta)\).
\emph{Parametric resonance} arises if the driving frequency \(\gamma\) matches certain resonance bands of the system.
\textbf{Effective Potential / Action Integral}
\begin{itemize}
\item From \(L = T - V\), one may define an \emph{effective potential} or use an action integral approach.
\item In many treatments, \(\int p\, d\phi = 2\pi\) (or a similar phase‐space integral) helps analyze periodic orbits and resonance conditions.
\end{itemize}

\vspace{0.5em}
\textbf{Polar / Angle Variables}
\[
\frac{\partial \sigma}{\partial S}
\;=\;
\frac{p(\theta)}{\sin\theta}
\,\frac{\partial p}{\partial \theta}.
\]
This type of relation can appear when converting to polar (or action‐angle) variables in a parametric oscillator system, relating momentum \(p(\theta)\) to the angular variable \(\theta\).

\end{document}


\end{document}
